\documentclass[12pt]{article}

\usepackage{setspace}
\onehalfspace

%\usepackage[utf8]{inputenc}
%\usepackage{t1enc}
%\usepackage[T1]{fontenc}

\usepackage{amsmath}

\usepackage{moodle}

\DeclareMathOperator{\Exp}{Exp}
\DeclareMathOperator{\Norm}{{\cal N}}
\DeclareMathOperator{\Binom}{Binom}
\DeclareMathOperator{\E}{E}
\DeclareMathOperator{\D}{D}
\renewcommand{\P}{\operatorname{P}}
%\DeclareMathOperator{\cov}{cov}
\DeclareMathOperator{\corr}{corr}

\DeclareMathOperator{\F}{F} 
\renewcommand{\d}{\operatorname{d\! }} %integral \! a nyero


\DeclareMathOperator{\R}{\mathbb{R}}
%\DeclareMathOperator{\Norm}{{\cal N}}
\DeclareMathOperator{\Unif}{{\cal U}}
\DeclareMathOperator{\Poi}{{Poisson}}


\DeclareMathOperator{\cov}{cov}





\begin{document}

\begin{quiz}{gyakorló}

\begin{multi}{eofv8}
Bármely $\xi$-re és $a<b$-re $\P(a< \xi <b)=F_{\xi}(b)-F_{\xi}(a)$.
\item igaz
\item* hamis
\end{multi}
\begin{multi}{impdvv5}
Azt mondjuk, hogy $\xi\sim \Binom(n,p)$, ha 
$$
\P(\xi=k)=\binom{n}{k}p^k(1-p)^{n-k}\ \ \ k=0\ldots n
$$
valamilyen $n>0$ és $p\in[0,1]$ esetén.
\item* igaz
\item hamis
\end{multi}
\begin{multi}{sfv3}
Azt mondjuk, hogy egy $f:\R\to \R$ sűrűségfüggvény 
$$
\int_{-\infty}^{+\infty}f(x)\d x\ge 0
$$
\item igaz
\item* hamis
\end{multi}
\begin{multi}{dvv10}
Egy $\xi$ valségi változó az $x_1,\ldots x_n$ értékeket veheti fel. 
Ekkor a szórásnégyzete:
$$
\D^2(\xi)=\E(\xi^2)-\E(\xi)^2
$$
\item* igaz
\item hamis
\end{multi}
\begin{multi}{impdvv6}
Azt mondjuk, hogy $\xi\sim \Binom(n,p)$, ha 
$$
\P(\xi=k)=\binom{n}{k}p^k(1-p)\ \ \ (k=0,1\ldots n)
$$
valamilyen $n>0$ és $p\in[0,1]$ esetén.
\item igaz
\item* hamis
\end{multi}
\begin{multi}{esem6}
Ha $\P(A)=1$ akkor $A$ mindig bekövetkezik.
\item igaz
\item* hamis
\end{multi}
\begin{multi}{impfvv11}
Ha $\xi\sim \Unif(a,b)$ akkor a sűrűségfüggvénye 
$$
f_{\xi}(x)=
\begin{cases}
\frac{1}{b-a} & \ \ a<x<b \\
0 & \text{máskor} 
\end{cases}
$$
\item* igaz
\item hamis
\end{multi}
\begin{multi}{val4}
A valószínűség monoton, azaz $\P(A)\le \P(B)$, ha $A\subseteq B$.
\item* igaz
\item hamis
\end{multi}
\begin{multi}{sfv1}
Azt mondjuk, hogy egy $f:\R\to \R$ sűrűségfüggvény, ha 
$$
f\ge 0 \ \ \text{és}\ \ \int_{-\infty}^{+\infty}f(x)\d x=1
$$
\item* igaz
\item hamis
\end{multi}
\begin{multi}{felt5}
$A_1,\ldots,A_n$ teljes eseményrendszer, ha $\sum A_k=\Omega$ és a tagok páronként kizáróak.
\item* igaz
\item hamis
\end{multi}
\begin{multi}{impfvv9}
Ha $\xi\sim \Exp(\lambda)$ akkor a sűrűségfüggvénye 
$$
f_{\xi}(x)=
\begin{cases}
\lambda e^{-\lambda x}& \ \ x>0\\
0 & \text{máskor} 
\end{cases}
$$
\item* igaz
\item hamis
\end{multi}
\begin{multi}{felt1}
$A$-nak a $B$-re vonatkozó feltételes valsége $\P(A|B)=\frac{\P(A)}{\P(B)}$.
\item igaz
\item* hamis
\end{multi}
\begin{multi}{flenvv9}
Ha a $\xi$ és $\eta$ valségi változóknak létezik a korrelációja, akkor $-1\le \corr(\xi,\eta)\le 1$.
\item* igaz
\item hamis
\end{multi}
\begin{multi}{esem4}
Az $\overline{A} + \overline{B}$ maga után vonja az $\overline{A\cdot B}$ eseményt.
\item* igaz
\item hamis
\end{multi}
\begin{multi}{flenvv1}
Azt mondjuk, hogy $\xi$ és $\eta$ diszkrét valségi változók függetlenek, ha
$$
\P(\xi=x,\eta=y)=\P(\xi=x)\P(\eta=y)
$$
bármely $x,y\in \R$-re teljesül.
\item* igaz
\item hamis
\end{multi}
\begin{multi}{eofv7}
Bármely $\xi$-re és $a<b$-re $\P(a\le \xi <b)=F_{\xi}(b)-F_{\xi}(a)$.
\item* igaz
\item hamis
\end{multi}
\begin{multi}{dvv7}
Egy $\xi$ valségi változó az $x_1,\ldots x_n$ értékeket veheti fel. Ekkor a szórásnégyzete:
$$
\D^2(\xi)=\sum_{k=1}^n \P(\xi=x_k)(x_k-\E(\xi))^2
$$
\item* igaz
\item hamis
\end{multi}
\begin{multi}{felt2}
$A$-nak a $B$-re vonatkozó feltételes valsége $\P(A|B)=\frac{\P(A)}{\P(B)}$, ha $A$ maga után vonja $B$-t.
\item* igaz
\item hamis
\end{multi}
\begin{multi}{dvv6}
Egy $\xi$ valségi változó az $x_1,\ldots x_n$ értékeket veheti fel. Ekkor a szórása:
$$
\D(\xi)=\sum_{k=1}^n \P(\xi=x_k)(x_k-\E(\xi))^2
$$
\item igaz
\item* hamis
\end{multi}
\begin{multi}{eofv6}
Diszkrét valségi változók eloszlásfüggvénye lépcsős függvény.
\item* igaz
\item hamis
\end{multi}
\begin{multi}{val5}
Bármely esemény valsége kiszámolható a $\frac{\text{kedvező}}{\text{összes}}$ képlettel.
\item igaz
\item* hamis
\end{multi}
\begin{multi}{dvv1}
Egy $p_1,\ldots, p_n$ vektor pont akkor valségeloszlás, ha $\sum_{k=1}^n p_k=1$.
\item igaz
\item* hamis
\end{multi}
\begin{multi}{impfvv2}
Egy $\xi\sim \Unif(a,b)$ várható értéke: $\frac{b-a}{2}$
\item igaz
\item* hamis
\end{multi}
\begin{multi}{sfv11}
Ha $\xi$-nek $F$ illetve $f$ az eloszlásfüggvény illetve sűrűségfüggvénye, akkor bármely $a<b$-re:
$$
F(b)-F(a)=\int_{a}^{b}f(x)\d x
$$
\item* igaz
\item hamis
\end{multi}
\begin{multi}{esem7}
Az $A$ és $B$ események pontosan akkor függetlenek, ha $\P(AB)=\P(A)\P(B)$.
\item* igaz
\item hamis
\end{multi}
\end{quiz}

\end{document}

